% Generated draft: requires mathematical and source review.
\documentclass[11pt]{article}
\usepackage[margin=25mm]{geometry}
\usepackage{fontspec}
\setmainfont{Latin Modern Roman}
\usepackage{amsmath,amssymb,mathtools,unicode-math}
\setmathfont{Latin Modern Math}
\usepackage{enumitem,needspace,xurl,hyperref}
\setlength{\parindent}{0pt}
\setlength{\parskip}{5pt}
\setlength{\emergencystretch}{3em}
\newcommand{\N}{\mathbb{N}}
\newcommand{\Z}{\mathbb{Z}}
\newcommand{\Q}{\mathbb{Q}}
\newcommand{\R}{\mathbb{R}}
\newcommand{\C}{\mathbb{C}}
\newcommand{\F}{\mathbb{F}}
\newcommand{\A}{\mathbb{A}}
\newcommand{\PP}{\mathbb{P}}
\newcommand{\E}{\mathbb{E}}
\newcommand{\eps}{\varepsilon}
\newcommand{\dd}{\,\mathrm{d}}
\begin{document}
\section*{UnsolvedMath ID: 30000870}
\textbf{Source problem code:} catalogue entry
\subsection{Definitions and mathematical statement}
The principal source question is: for every irreducible Hodge structure $V$ in the intended category, is
\[
Z\bigl(\operatorname{Hg}(V)\bigr)=\mathbb G_m?
\]
Here $Z(G)$ denotes the algebraic center of an algebraic group $G$, and $\mathbb G_m$ is the multiplicative algebraic group. Both the category of Hodge structures and the meaning of this equality require clarification.

For orientation, under the conventional rational pure-Hodge interpretation, $V$ is a finite-dimensional $\Q$-vector space with a decomposition $V_{\C}=\bigoplus_{p+q=w}V^{p,q}$ satisfying $\overline{V^{p,q}}=V^{q,p}$; irreducibility means absence of nonzero proper rational Hodge substructures. A common convention defines the Mumford--Tate group from the Hodge homomorphism $h:\operatorname{Res}_{\C/\R}\mathbb G_m\to\operatorname{GL}(V_{\R})$, and the Hodge group from its restriction to the norm-one circle. These are different groups: a scalar weight torus present in a Mumford--Tate formulation cannot simply be transferred to the Hodge-group question. The source alternates between a question written using $\operatorname{Hg}(V)$ and an assessment about Mumford--Tate groups. It does not specify whether the full center or its identity component is intended, whether $\mathbb G_m$ means the scalar subgroup or an abstract algebraic-group isomorphism, or whether polarizability and weight restrictions are imposed. Consequently the displayed assertion is retained as the literal question, without asserting a universal affirmative answer or replacing its group.

The original statement also asks whether there exists an abelian variety $A$ whose associated Weil-group representation $\rho$ is not equivalent to a direct sum of monomial representations. Conventionally, a monomial representation is induced from a one-dimensional character of a suitable subgroup; the relevant Weil group, coefficient field, induction category and construction of $\rho$ are not supplied here. The source expressly reports an affirmative example for this second question, but supplies no example or construction in the record. Thus this is an attributed answered question, not an additional unresolved conjecture in this draft.

The supplied triage dated 2026-08-21 labels the record open and broadly calls the center and monomiality questions unresolved. Its monomiality assessment conflicts with the original statement's reported affirmative example; both qualifications are retained.
\subsection{Short English statement}
The main question asks whether every irreducible Hodge structure has Hodge-group center $\mathbb G_m$, but the intended group and center conventions are unclear. A second source question asks for an abelian variety with a nonmonomial Weil representation and is explicitly reported there to have an affirmative example.
\subsection{Sources}
\begin{itemize}
\item Local UnsolvedMath registry, permanent ID 30000870, source code catalogue entry; supplied complete source record, including its longer original statement and dated literature assessment.
\item Modulformen (2008), Oberwolfach Reports, DOI 10.4171/owr/2007/51, pp. 3013--3072: \url{https://doi.org/10.4171/owr/2007/51}. The record gives proposed year 2007, no named proposer, and a corrected-verified editorial-cleanup label.
\end{itemize}
\end{document}
